\subsection{The closure of a linear variety}

Recall (A, II, § 9.3) that in a vector space E over a division ring K, a non-empty affine linear variety (called « linear variety » when this can cause no confusion) is the image under a translation of a vector subspace of E.

PROPOSITION 1. --- In a topological vector space E, the closure of a linear variety is a linear variety.

Since every translation is a homeomorphism of E, it is sufficient to demonstrate the proposition for a vector subspace M of E, and in this case, the proposition has been proved in I, p.~4.

COROLLARY. --- In a topological vector space E, every hyperplane is either closed or everywhere dense.

In fact, the closure of a homogeneous hyperplane H can only be H or the whole space E, since it is a vector subspace containing H (prop. 1).

A hyperplane H is closed in E if, and only if, CH contains an interior point.

The vector subspace M generated by a set A, in a topological vector space E, is the set of linear combinations of points of A (A, II, § 1.7, prop. 9); the closure of M in E is, by prop. 1, the smallest closed vector subspace containing A; we say that this is the closed vector subspace generated by A.

DEFINITION 1. --- A set A, in a topological vector space E, is total if, and only if, the closed vector subspace generated by A coincides with E (i.e.~the set of linear combinations of elements of A is everywhere dense).

Examples. --- 1) In the normed space \(\mathcal{C}(\mathrm{I};\mathbf{C})\) (on the field \(\mathbf{C}\) ) of functions, continuous on \(\mathrm{I} = [0,1]\) , with values in \(\mathbf{C}\) , the restrictions to \(\mathrm{I}\) of the functions \(x^n\) ( \(n\in \mathbf{N}\) ) form a total set, by the Weierstrass-Stone theorem (GT, X, § 4.2, th. 3). Similarly, the restrictions to \(\mathrm{I}\) of the functions \(e^{2n\pi ix}\) ( \(n\in \mathbf{Z}\) ) form a total set (GT, X, § 4.4, prop. 8), in the subspace \(\mathbf{P}\) of \(\mathcal{C}(\mathrm{I},\mathbf{C})\) formed of functions such that \(f(0) = f(1)\) .

\begin{enumerate}
\def\labelenumi{\arabic{enumi})}
\setcounter{enumi}{1}
\tightlist
\item
  Every absorbent set in a topological vector space E over a non-discrete valued division ring (and in particular every neighbourhood of 0 in E) is a total set since it generates E (I, p.~7). Thus a linear variety that is not dense in E is necessarily a nowhere dense set in E (GT, IX, § 5.1) since its closure cannot contain an interior point.
\end{enumerate}

DEFINITION 2. --- A family \((a_{\iota})_{\iota\in\mathbf{I}}\) of points of a topological vector space E is called topologically independent if for any \(\kappa\in I\) , the closed vector subspace generated by the \(a_{\iota}\) , with \(\iota\neq\kappa\) , does not contain \(a_{\kappa}\) .

Example. --- 3) In the normed space \(\mathcal{C}(\mathrm{I};\mathbf{C})\) of continuous functions defined over \(\mathrm{I} = [0,1]\) , the restrictions to \(\mathrm{I}\) of the functions \(e^{2n\pi ix}(n\in \mathbf{Z})\) form a topologically independent family. If \(f(x)\) is the linear combination \(\sum_{k\neq n}c_k e^{2k\pi ix}\) (where all but finitely many of the \(c_k\) are zero) then

\[
\int_ {0} ^ {1} \left| e ^ {2 n \pi i x} - f (x) \right| ^ {2} d x = 1 + \sum_ {k \neq n} | c _ {k} | ^ {2} \geqslant 1
\]

and, a fortiori, by the mean value theorem

\[
\sup _ {x \in \mathbf {I}} \left| e ^ {2 n \pi i x} - f (x) \right| \geqslant 1
\]

which shows that \(e^{2\pi inx}\) does not belong to the closed vector subspace of \(\mathcal{C}(\mathrm{I};\mathbf{C})\) generated by \(e^{2k\pi ix}, k \neq n\) .

The set of elements of a topologically independent family is called a topologically independent set of E. Every subset of a topologically independent subset is topologically independent; every subset consisting of a single point \(x \neq 0\) is topologically independent if E is a Hausdorff space.

A topologically independent family is independent (in the algebraic sense; cf.~A, II, § 7.1, Remark), but the converse is incorrect.

Example. --- 4) In the normed space \(\mathcal{C}(\mathrm{I};\mathbf{C})\) of functions that are continuous over \(\mathrm{I} = [0,1]\) , the restriction to \(\mathrm{I}\) of the functions \(x^n\) ( \(n\in \mathbf{N}\) ) form an algebraically independent family. But there exists a sequence of polynomials ( \(p_n\) ) such that \(p_n(x^2)\) converges uniformly to \(x\) in \(\mathrm{I}\) (GT, X, § 4.2, lemma 2) which shows that \(x\) belongs to the closed vector subspace of \(\mathcal{C}(\mathrm{I};\mathbf{C})\) generated by the functions \(x^{2n}\) ( \(n\in \mathbf{N}\) ).

Remarks. --- 1) The family of topologically independent sets of a topological vector space is not necessarily inductive for the relation of inclusion (I, p.~25, exerc. 2); this situation is thus different to that for algebraically independent sets. Moreover there does not necessarily exist in E a maximal topologically independent subset (I, p.~25, exerc. 4), thus there does not necessarily exist a subset that is both total and topologically independent.

\begin{enumerate}
\def\labelenumi{\arabic{enumi})}
\setcounter{enumi}{1}
\tightlist
\item
  Let \(\mathbf{M}\) be a closed vector subspace of \(\mathrm{E}\) and \((\dot{a}_{\mathrm{t}})_{\mathrm{t}\in \mathrm{I}}\) a topologically independent family in the quotient space \(\mathrm{E / M}\) . If \(a_{\mathrm{t}}\) is any element of the class \(\dot{a}_{\mathrm{t}}\) , then from def. 2, and the fact that the canonical mapping of \(\mathrm{E}\) on \(\mathrm{E / M}\) is continuous, it follows that the family \((a_{\mathrm{t}})_{\mathrm{t}\in \mathrm{I}}\) is topologically independent. But note that if \(\mathbf{N}\) is the closed vector subspace generated by the \(a_{\mathrm{t}}\) it can happen that \(\mathbf{M} \cap \mathbf{N} \neq \{0\}\) (I, p.~25, exerc. 2), and hence the sum \(\mathbf{M} + \mathbf{N}\) is not necessarily direct in the algebraic sense (nor a fortiori in the topological sense).
\end{enumerate}

